\begin{helper}
For every \(\varepsilon_2>0\), there is \(n_0\) such that for every
\(n>n_0\), if
\[
L=\log n,\quad
m=(\noderef{TwoBiteNaturalReducedVertexCount}(n):\mathbb R),\quad
p=\frac12\sqrt{\frac Ln},\quad
t_1=\frac{\sqrt{nL}}{\log L},
\]
then the hypotheses \(0<n\), \(0<L\), \(0<\log L\), and
\[
2pm\le\frac{\sqrt n}{L^{3/2}}
\]
imply
\[
2pm\le\frac{\varepsilon_2}{100}t_1.
\]
\end{helper}
\begin{proof}
Instantiate \(\noderef{ParameterHierarchyT9LogThreshold}\) with
\(\varepsilon_2/5\), which is positive.  After increasing \(n_0\), for every
\(n>n_0\) we have
\[
\frac{2\log L}{L^2}\le \frac{\varepsilon_2/5}{10}
=\frac{\varepsilon_2}{50}.
\]
Dividing by \(2\) gives
\[
\frac{\log L}{L^2}\le\frac{\varepsilon_2}{100}.
\]
Now apply \(\noderef{ParameterHierarchyT16Algebra}\) with this threshold
bound, together with the assumed positivity of \(n\), \(L\), and \(\log L\),
and with the assumed mass estimate
\[
2pm\le\frac{\sqrt n}{L^{3/2}}.
\]
The conclusion is exactly
\[
2pm\le\frac{\varepsilon_2}{100}t_1.
\]
\end{proof}
